% ============================================================
%  ch00  开篇、路线图、记号与约定、致谢
%  原文对应：ICM-arXiv.tex 486--526 行
% ============================================================

经典朗兰兹纲领由 Langlands 于 1960 年代创立\cite{Langlands-to-Weil}，在表示论的框架下系统地研究互反律。粗略地说，它预测数论与表示论之间存在下述深刻联系。

\[
\begin{tikzcd}[cells={nodes={draw, inner sep=5pt}}, column sep=4.4cm, row sep=1.7cm]
  \text{Galois 表示} \arrow[r, <->, "{\scriptsize\text{互反律}}" above] \arrow[d, <->] &
  \text{自守表示} \arrow[d, <->] \\
  \text{算术数据} \arrow[r, <->, "{\scriptsize\text{Satake 同构}}" above] &
  \text{谱数据}
\end{tikzcd}
\]

这一对应的一类特例称为志村--Taniyama--Weil 猜想，它蕴含 Fermat 大定理（见 \cite{Wiles-Fermat}）。\zy{原文写作 ``Shimura-Tanniyama-Weil conjecture''，其中 Tanniyama 为 Taniyama 之误植。此处按通行拼法作 Taniyama，原样登记于 errata 表。}

几何朗兰兹纲领\cite{Laumon-Geometric-Langlands} 由 Drinfeld 与 Laumon 发起，其来源是 Drinfeld 处理函数域上 $\GL_2$ 整体朗兰兹对应的途径\cite{Drinfeld-GL2} 的推广。在几何理论中，基本研究对象从约化群 $G$ 的自守形式空间，转移到了代数曲线上 $G$-丛模空间上的层范畴。

长期以来，几何朗兰兹的发展受到经典朗兰兹中问题与技巧的启发，而另一个重要灵感来源是量子物理。其基本哲学称为范畴化／几何化。然而近些年来，几何理论反过来在经典朗兰兹纲领以及若干相关算术问题上找到了丰硕的应用。传统上，人们用 Grothendieck 的层--函数字典把几何理论中研究的层"去范畴化"，得到算术理论中研究的函数。前面提到的 Drinfeld 处理 $\GL_2$ 朗兰兹对应的工作就用了这一点。另一个著名例子是 Ngô 对基本引理的证明（\cite{Ngo-Fundamental-Lemma}）。近年来出现了从几何理论通往算术理论的另一条通道，同样借助迹构造，但性质不同，且与物理中的想法密切相关。V. Lafforgue 关于函数域上整体朗兰兹对应的工作（\cite{Lafforgue-Automorphic-to-Galois}）本质上（但并不显式地）用到了这一想法。

本文综述几何朗兰兹的（一小部分）进展，并讨论经典朗兰兹的一些近期新视角——它们由几何理论启发而来，又反过来导致若干具体算术问题的解决。下图可视为一张路线图。

\[
\begin{tikzpicture}[font=\small]
  \node[draw, inner sep=4pt, minimum width=4.8cm] (n1) {数论、算术几何、调和分析};
  \node[draw, inner sep=7pt, minimum width=4.8cm, below=11mm of n1] (n2) {\large 算术朗兰兹纲领};
  \node[draw, inner sep=7pt, minimum width=4.8cm, below=40mm of n2] (n3) {\large 几何朗兰兹纲领};
  \node[draw, inner sep=4pt, minimum width=4.8cm, below=11mm of n3] (n4) {表示论、几何、量子物理};
  \draw[<->] (n1) -- (n2);
  \draw[<->] (n3) -- (n4);
  \draw[double, -{Implies}, thick] (n2.east) .. controls +(2.5cm,0) and +(2.5cm,0) .. (n3.east)
    node[midway, right=4pt] {范畴化／几何化};
  \draw[double, -{Implies}, thick] (n3.west) .. controls +(-2.5cm,0) and +(-2.5cm,0) .. (n2.west)
    node[midway, left=4pt] {去范畴化／迹};
\end{tikzpicture}
\]

\subsection*{记号与约定}
\addcontentsline{toc}{subsection}{记号与约定}

本文通篇使用以下记号。

设 $F$ 为域，$\widetilde F/F$ 为 Galois 扩张，以 $\Ga_{\widetilde F/F}$ 记该扩张的 Galois 群。令 $\Ga_F=\Ga_{\overline F/F}$，其中 $\overline F$ 是 $F$ 的一个可分闭包。文中 $F$ 常常取局部域或整体域；此时以 $W_F$ 记 $F$ 的 Weil 群。以 $\cycl$ 记分圆特征。

设 $F$ 为域，$A$ 为 $F$ 上的乘法型群。以 $\xch(A)=\Hom(A_{\overline F},\bG_m)$ 记其特征群，以 $\xcoch(A)=\Hom(\bG_m,A_{\overline F})$ 记其余特征群。

设 $\ell$ 为素数。令 $\La$ 为 $\bF_\ell$、$\bZ_\ell$、$\bQ_\ell$，或这类环的有限（平坦）扩张。它将充当本文层理论的系数环。

\subsection*{致谢}
\addcontentsline{toc}{subsection}{致谢}

作者感谢他所有的合作者：没有他们，本综述中报告的许多工作都无从谈起。此外，他还感谢 Edward Frenkel、Dennis Gaitsgory、何旭华、Michael Rapoport、Peter Scholze 多年来在朗兰兹纲领各部分上对他的教导与讨论。本项工作部分受到 NSF（协议号 DMS-1902239）以及 Simons fellowship 的资助。
